% Options for packages loaded elsewhere
\PassOptionsToPackage{unicode}{hyperref}
\PassOptionsToPackage{hyphens}{url}
\documentclass[
]{article}
\usepackage{xcolor}
\usepackage[margin=1in]{geometry}
\usepackage{amsmath,amssymb}
\setcounter{secnumdepth}{-\maxdimen} % remove section numbering
\usepackage{iftex}
\ifPDFTeX
  \usepackage[T1]{fontenc}
  \usepackage[utf8]{inputenc}
  \usepackage{textcomp} % provide euro and other symbols
\else % if luatex or xetex
  \usepackage{unicode-math} % this also loads fontspec
  \defaultfontfeatures{Scale=MatchLowercase}
  \defaultfontfeatures[\rmfamily]{Ligatures=TeX,Scale=1}
\fi
\usepackage{lmodern}
\ifPDFTeX\else
  % xetex/luatex font selection
\fi
% Use upquote if available, for straight quotes in verbatim environments
\IfFileExists{upquote.sty}{\usepackage{upquote}}{}
\IfFileExists{microtype.sty}{% use microtype if available
  \usepackage[]{microtype}
  \UseMicrotypeSet[protrusion]{basicmath} % disable protrusion for tt fonts
}{}
\makeatletter
\@ifundefined{KOMAClassName}{% if non-KOMA class
  \IfFileExists{parskip.sty}{%
    \usepackage{parskip}
  }{% else
    \setlength{\parindent}{0pt}
    \setlength{\parskip}{6pt plus 2pt minus 1pt}}
}{% if KOMA class
  \KOMAoptions{parskip=half}}
\makeatother
\usepackage{longtable,booktabs,array}
\usepackage{calc} % for calculating minipage widths
% Correct order of tables after \paragraph or \subparagraph
\usepackage{etoolbox}
\makeatletter
\patchcmd\longtable{\par}{\if@noskipsec\mbox{}\fi\par}{}{}
\makeatother
% Allow footnotes in longtable head/foot
\IfFileExists{footnotehyper.sty}{\usepackage{footnotehyper}}{\usepackage{footnote}}
\makesavenoteenv{longtable}
\usepackage{graphicx}
\makeatletter
\newsavebox\pandoc@box
\newcommand*\pandocbounded[1]{% scales image to fit in text height/width
  \sbox\pandoc@box{#1}%
  \Gscale@div\@tempa{\textheight}{\dimexpr\ht\pandoc@box+\dp\pandoc@box\relax}%
  \Gscale@div\@tempb{\linewidth}{\wd\pandoc@box}%
  \ifdim\@tempb\p@<\@tempa\p@\let\@tempa\@tempb\fi% select the smaller of both
  \ifdim\@tempa\p@<\p@\scalebox{\@tempa}{\usebox\pandoc@box}%
  \else\usebox{\pandoc@box}%
  \fi%
}
% Set default figure placement to htbp
\def\fps@figure{htbp}
\makeatother
\setlength{\emergencystretch}{3em} % prevent overfull lines
\providecommand{\tightlist}{%
  \setlength{\itemsep}{0pt}\setlength{\parskip}{0pt}}
\usepackage{bookmark}
\IfFileExists{xurl.sty}{\usepackage{xurl}}{} % add URL line breaks if available
\urlstyle{same}
\hypersetup{
  hidelinks,
  pdfcreator={LaTeX via pandoc}}

\author{}
\date{\vspace{-2.5em}}

\begin{document}

\section{Relatório de Revisão Numérica
Cruzada}\label{relatuxf3rio-de-revisuxe3o-numuxe9rica-cruzada}

\textbf{Dissertação de Mestrado --- Letícia Schimidt Arruda (FMUSP)}\\
\textbf{Revisor:} Sub-agente Sonnet (revisão automática)\\
\textbf{Data:} 19 de maio de 2026\\
\textbf{Arquivos revisados:} \texttt{index.qmd},
\texttt{manuscript/tabelas\_dissertacao.docx},
\texttt{manuscript/resultados\_redacao.docx},
\texttt{results/tab\_obj2\_*.csv},
\texttt{results/tabelas\_dissertacao/tab*.csv},
\texttt{results/Relatorio\ Analise\ SPSS.md},
\texttt{manuscript/Dissertação\ Maio\_26.docx}

\begin{center}\rule{0.5\linewidth}{0.5pt}\end{center}

\subsection{Sumário Executivo}\label{sumuxe1rio-executivo}

\begin{itemize}
\tightlist
\item
  \textbf{OK:} Tabela 10 (Modelo 4) está numericamente perfeita ---
  todos os OR, IC 95\% e p-valores batem exatamente com o Relatorio
  Analise SPSS.md (linhas 123--135). Forest plot
  (\texttt{fig\_obj5\_forest\_plot\_modelo4.png}) foi gerado e
  corresponde aos mesmos valores.
\item
  \textbf{OK:} Tabelas 5 e 6 são internamente consistentes --- para cada
  célula, \texttt{n\_cesárea\ /\ n\_total} coincide com a taxa
  percentual da Tabela 6 (verificado em todos os 12 grupos com dados
  disponíveis).
\item
  \textbf{OK:} Todos os números de comorbidades (§4.2) e de regressão
  logística (§4.5) na redação têm origem rastreável e conferem com o
  Relatorio SPSS.
\item
  \textbf{CRÍTICO:} Existem quatro valores distintos de N total
  circulando nos documentos (6.646, 6.663, 6.830 e 8.300/8.724). Nenhuma
  tabela usa exatamente o mesmo N, e os cinco placeholders da redação
  ainda não foram preenchidos.
\item
  \textbf{IMPORTANTE:} CSVs derivados (\texttt{tab05/tab06\_*.csv}) têm
  valores diferentes dos CSVs primários (\texttt{tab\_obj2\_*.csv}) em
  quatro células --- as Tabelas 5--7 do docx foram construídas
  corretamente a partir dos primários, mas os derivados contêm um bug de
  recomputação.
\item
  \textbf{IMPORTANTE:} Figuras 5 e 6 (indicações de cesárea e fórcipe)
  são citadas na redação §4.4 mas os arquivos PNG não existem em
  \texttt{results/figures/}. Figura 7 (forest plot) existe mas não é
  referenciada na redação.
\end{itemize}

\begin{center}\rule{0.5\linewidth}{0.5pt}\end{center}

\subsection{1. Inconsistências
Encontradas}\label{inconsistuxeancias-encontradas}

\subsubsection{1.1 N total da amostra --- quatro valores em circulação
(CRÍTICO)}\label{n-total-da-amostra-quatro-valores-em-circulauxe7uxe3o-cruxedtico}

\textbf{Arquivos:} \texttt{Dissertação\ Maio\_26.docx},
\texttt{tabelas\_dissertacao.docx},
\texttt{tab\_obj2\_vias\_parto\_geral.csv},
\texttt{tab01\_sociodemografia.csv}

\begin{longtable}[]{@{}
  >{\raggedright\arraybackslash}p{(\linewidth - 8\tabcolsep) * \real{0.2000}}
  >{\raggedright\arraybackslash}p{(\linewidth - 8\tabcolsep) * \real{0.2000}}
  >{\raggedright\arraybackslash}p{(\linewidth - 8\tabcolsep) * \real{0.2000}}
  >{\raggedright\arraybackslash}p{(\linewidth - 8\tabcolsep) * \real{0.2000}}
  >{\raggedright\arraybackslash}p{(\linewidth - 8\tabcolsep) * \real{0.2000}}@{}}
\toprule\noalign{}
\begin{minipage}[b]{\linewidth}\raggedright
Fonte
\end{minipage} & \begin{minipage}[b]{\linewidth}\raggedright
Precoces
\end{minipage} & \begin{minipage}[b]{\linewidth}\raggedright
Tardias
\end{minipage} & \begin{minipage}[b]{\linewidth}\raggedright
Adultas
\end{minipage} & \begin{minipage}[b]{\linewidth}\raggedright
Total
\end{minipage} \\
\midrule\noalign{}
\endhead
\bottomrule\noalign{}
\endlastfoot
R pipeline --- \texttt{tab\_obj2\_vias\_parto\_geral.csv} (primário) &
538 & 829 & 5.279 & \textbf{6.646} \\
Tabelas 4--7 no \texttt{tabelas\_dissertacao.docx} (rodapé Tab4) & 538 &
830 & 5.295 & \textbf{6.663} \\
SPSS --- cabeçalho da Tabela 1 no \texttt{tabelas\_dissertacao.docx} &
551 & 856 & 5.423 & \textbf{6.830} \\
Dissertação §154 (texto corrido) & --- & 1.348 adol & 6.952 adul &
\textbf{8.300} \\
Dissertação legenda Tabelas 1--2 & --- & 1.677 adol & 7.047 adul &
\textbf{8.724} \\
\end{longtable}

\textbf{Problema:} Há quatro valores de N distintos, cada um
correspondendo a um escopo diferente de filtragem: - \textbf{6.646} = R
com todos os critérios de elegibilidade aplicados (IG ≥ 22, Robson 1--10
exceto 8, idade 11--34, tipo\_parto preenchido). - \textbf{6.663} =
Tabelas 5--7 do docx: soma dos n por Robson a partir de
\texttt{tab\_obj2\_taxa\_cesarea.csv} --- 17 casos a mais que o
primário, sugerindo que a tabela de taxa e a de vias usaram cortes de
código diferentes no mesmo render. - \textbf{6.830} = SPSS: inclui os
184 casos que o R exclui por Robson ausente/inválido ou IG \textless{}
22. - \textbf{8.300} e \textbf{8.724} = dissertação original: inclui
gestantes ≥ 35 anos (n ≈ 1.469 no SPSS) e casos pré-filtro; os dois
valores são incompatíveis entre si (diferença de 424), indicando que
§154 e as legendas das Tabelas 1--2 da dissertação usam cortes
distintos.

\textbf{Correção sugerida:} Adotar \textbf{6.663} como N definitivo das
análises de via de parto (Tabelas 4--7), pois é o valor que soma os n de
todos os grupos de Robson no docx e é internamente consistente com a
soma das células das Tabelas 5--7. Adotar \textbf{6.830} (ou o n exato
que o SPSS usou) para as Tabelas 1--3 (sociodemografia, hábitos,
comorbidades), com nota clara de que o denominador difere por incluir
casos sem Robson/IG. Eliminar as referências a 8.300 e 8.724 no texto
corrido ou explicar explicitamente que esses são os totais brutos
pré-critérios. Os cinco placeholders ({[}N\_TOTAL{]}, {[}N\_ADOL{]},
{[}N\_PREC{]}, {[}N\_TARD{]}, {[}N\_ADUL{]}) devem ser preenchidos com
os valores da tabela de análise escolhida (ver Seção 3).

\begin{center}\rule{0.5\linewidth}{0.5pt}\end{center}

\subsubsection{1.2 Discrepância entre CSVs primários e CSVs derivados
(Tabelas
5--7)}\label{discrepuxe2ncia-entre-csvs-primuxe1rios-e-csvs-derivados-tabelas-57}

\textbf{Arquivos:} \texttt{tab\_obj2\_taxa\_cesarea.csv} vs
\texttt{tab06\_taxa\_cesarea\_robson.csv};
\texttt{tab\_obj2\_robson\_x\_parto\_x\_idade.csv} vs
\texttt{tab05\_robson\_faixa\_via.csv}

Quatro células diferem entre os CSVs primários (\texttt{tab\_obj2\_*}) e
os derivados (\texttt{tab0X\_*}):

\begin{longtable}[]{@{}
  >{\raggedright\arraybackslash}p{(\linewidth - 6\tabcolsep) * \real{0.2500}}
  >{\raggedright\arraybackslash}p{(\linewidth - 6\tabcolsep) * \real{0.2500}}
  >{\raggedright\arraybackslash}p{(\linewidth - 6\tabcolsep) * \real{0.2500}}
  >{\raggedright\arraybackslash}p{(\linewidth - 6\tabcolsep) * \real{0.2500}}@{}}
\toprule\noalign{}
\begin{minipage}[b]{\linewidth}\raggedright
Célula
\end{minipage} & \begin{minipage}[b]{\linewidth}\raggedright
Primário (\texttt{tab\_obj2\_taxa})
\end{minipage} & \begin{minipage}[b]{\linewidth}\raggedright
Derivado (\texttt{tab06})
\end{minipage} & \begin{minipage}[b]{\linewidth}\raggedright
Diferença
\end{minipage} \\
\midrule\noalign{}
\endhead
\bottomrule\noalign{}
\endlastfoot
Grupo 6 --- Adultas & n = 39, taxa = 94,9\% & n = 38, taxa = 94,7\% & −1
caso \\
Grupo 7 --- Adultas & n = 302, taxa = 90,4\% & n = 300, taxa = 90,7\% &
−2 casos \\
Grupo 10 --- Tardias & n = 116, taxa = 31,0\% & n = 115, taxa = 31,3\% &
−1 caso \\
Grupo 10 --- Adultas & n = 1.193, taxa = 62,0\% & n = 1.180, taxa =
62,5\% & −13 casos \\
\end{longtable}

\textbf{Problema:} As Tabelas 5--7 do \texttt{tabelas\_dissertacao.docx}
foram construídas corretamente a partir dos CSVs primários
(\texttt{tab\_obj2\_*}). Os CSVs derivados (\texttt{tab05/tab06\_*.csv})
foram recomputados em um momento diferente e produziram valores
levemente distintos --- provavelmente por diferença de seed de simulação
Monte Carlo, ordem de operações ou versão de pacote. A inconsistência
está nos derivados, não no docx.

\textbf{Correção sugerida:} Ignorar os CSVs derivados
(\texttt{tab05/tab06}) para a dissertação; usar exclusivamente os CSVs
primários (\texttt{tab\_obj2\_*}) como fonte de verdade para Tabelas
4--7. Deletar ou renomear os derivados para evitar confusão futura.

\begin{center}\rule{0.5\linewidth}{0.5pt}\end{center}

\subsubsection{1.3 Discrepância entre totais de Tabela 4 no docx e CSV
primário}\label{discrepuxe2ncia-entre-totais-de-tabela-4-no-docx-e-csv-primuxe1rio}

\textbf{Arquivo:} \texttt{tabelas\_dissertacao.docx} Tabela 4 vs
\texttt{tab\_obj2\_vias\_parto\_geral.csv}

\begin{longtable}[]{@{}lll@{}}
\toprule\noalign{}
Grupo & CSV primário (\texttt{tab\_obj2\_vias}) & Tab4 docx \\
\midrule\noalign{}
\endhead
\bottomrule\noalign{}
\endlastfoot
Tardias & n = 829 & n = 830 (+1) \\
Adultas & n = 5.279 & n = 5.295 (+16) \\
\textbf{Total} & \textbf{6.646} & \textbf{6.663} \\
\end{longtable}

\textbf{Problema:} A Tabela 4 no docx exibe totais diferentes dos totais
no CSV primário de vias de parto. Isso indica que dois CSVs produzidos
pelo mesmo \texttt{index.qmd} ---
\texttt{tab\_obj2\_vias\_parto\_geral.csv} e
\texttt{tab\_obj2\_taxa\_cesarea.csv} --- não são perfeitamente
concordantes. A soma dos n por Robson a partir de
\texttt{tab\_obj2\_taxa\_cesarea.csv} produz 6.663, enquanto
\texttt{tab\_obj2\_vias\_parto\_geral.csv} produz 6.646. Ambos deveriam
ser idênticos, pois derivam do mesmo \texttt{dados\_analise}.

\textbf{Correção sugerida:} Rerodar o \texttt{index.qmd} e verificar se
os dois CSVs convergem para o mesmo total. Se não convergirem,
identificar o chunk que produz cada um e garantir que ambos filtram a
mesma base. O N definitivo para as Tabelas 4--7 deve ser um único número
consistente.

\begin{center}\rule{0.5\linewidth}{0.5pt}\end{center}

\subsubsection{1.4 Dissertação original --- p-valor de
cor/raça}\label{dissertauxe7uxe3o-original-p-valor-de-corrauxe7a}

\textbf{Arquivo:} \texttt{Dissertação\ Maio\_26.docx} §160 vs
\texttt{resultados\_redacao.docx} §4.1.2 vs
\texttt{Relatorio\ Analise\ SPSS.md}

\begin{longtable}[]{@{}lll@{}}
\toprule\noalign{}
Fonte & p-valor & Conclusão \\
\midrule\noalign{}
\endhead
\bottomrule\noalign{}
\endlastfoot
Dissertação original §160 & p = \textbf{0,06} & não significativo \\
Relatório SPSS & p = \textbf{0,337} & não significativo \\
Nova redação (§4.1.2) & p = \textbf{0,337} & não significativo \\
Tabela 1 docx & p = \textbf{0,337} & não significativo \\
\end{longtable}

\textbf{Problema:} O texto original da dissertação (§160) cita p = 0,06,
que diverge do valor do SPSS (0,337). Ambos apontam ausência de
significância estatística, de modo que a conclusão clínica não muda. A
nova redação e a Tabela 1 já foram corrigidas para 0,337. O problema
residual é que §160 da dissertação original ainda contém o valor errado.

\textbf{Correção sugerida:} Na revisão do docx principal, substituir ``p
valor neste caso é de 0,06'' por ``p = 0,337'' em §160. O p definitivo é
0,337, proveniente do SPSS com a amostra conforme os critérios de
elegibilidade.

\begin{center}\rule{0.5\linewidth}{0.5pt}\end{center}

\subsubsection{1.5 Dissertação original --- incompatibilidade entre §154
e legendas das Tabelas
1--2}\label{dissertauxe7uxe3o-original-incompatibilidade-entre-154-e-legendas-das-tabelas-12}

\textbf{Arquivo:} \texttt{Dissertação\ Maio\_26.docx}

\begin{itemize}
\tightlist
\item
  §154 (texto): \emph{``8.300 participantes, sendo 1.348 parturientes
  adolescentes e 6.952 adultas''}
\item
  Legenda Tabela 1: \emph{``incluindo 1.677 parturientes adolescentes
  (11 a 19 anos) e 7.047 adultas''}
\item
  Legenda Tabela 2: \emph{``incluindo 1.677 adolescentes (de 11 a 19
  anos) e 7.047 adultas''}
\end{itemize}

\textbf{Problema:} §154 e as legendas das Tabelas 1--2 referem-se ao
mesmo estudo mas apresentam totais incompatíveis: 1.348 adol vs 1.677
adol (diferença de 329) e 6.952 adul vs 7.047 adul (diferença de 95).
Nenhum desses valores corresponde aos números do pipeline do R
(538+829=1.367 adol; 5.279--5.295 adul) nem ao SPSS filtrado
(551+856=1.407 adol; 5.423 adul).

\textbf{Correção sugerida:} Eliminar ambas as versões do texto original;
substituir pelo N definitivo pós-reconciliação (ver Seção 1.1 e Seção
3). Enquanto os placeholders não forem preenchidos, manter apenas a nova
redação (\texttt{resultados\_redacao.docx}) e não colar nenhum parágrafo
da versão original de §154 e legendas.

\begin{center}\rule{0.5\linewidth}{0.5pt}\end{center}

\subsubsection{1.6 Redação §4.1.2 --- percentual de cor/raça pode não
representar todas as
adolescentes}\label{redauxe7uxe3o-4.1.2-percentual-de-corrauxe7a-pode-nuxe3o-representar-todas-as-adolescentes}

\textbf{Arquivo:} \texttt{resultados\_redacao.docx} §4.1.2

\textbf{Trecho:} \emph{``56,3\% das adolescentes foram classificadas
como brancas''}

\textbf{Problema:} A Tabela 1 do docx mostra precoces = 56,6\% brancas e
tardias = 56,3\% brancas. O valor 56,3\% coincide com as tardias
isoladamente. O SPSS calcula a proporção com adolescentes unificadas
(precoces + tardias combinadas), e pode produzir um valor diferente de
56,3\% dependendo do peso de cada subgrupo. A redação não indica de qual
fonte veio esse 56,3\%.

\textbf{Correção sugerida:} Verificar no arquivo
\texttt{Resultados\_finais.ods} o percentual combinado de brancas para o
grupo adolescente (11--19 anos). Se diferir de 56,3\%, corrigir na
redação. Se for 56,3\%, acrescentar nota de que é o valor agregado do
SPSS, diferente dos 56,6\% (precoces) e 56,3\% (tardias) da Tabela 1.

\begin{center}\rule{0.5\linewidth}{0.5pt}\end{center}

\subsubsection{1.7 Tabela 4 docx vs SPSS --- diferença de percentuais de
vias de
parto}\label{tabela-4-docx-vs-spss-diferenuxe7a-de-percentuais-de-vias-de-parto}

\textbf{Arquivos:} \texttt{tabelas\_dissertacao.docx} Tabela 4 (fonte R)
vs \texttt{Relatorio\ Analise\ SPSS.md} §4 (tabela de via de parto)

\begin{longtable}[]{@{}
  >{\raggedright\arraybackslash}p{(\linewidth - 12\tabcolsep) * \real{0.1429}}
  >{\raggedright\arraybackslash}p{(\linewidth - 12\tabcolsep) * \real{0.1429}}
  >{\raggedright\arraybackslash}p{(\linewidth - 12\tabcolsep) * \real{0.1429}}
  >{\raggedright\arraybackslash}p{(\linewidth - 12\tabcolsep) * \real{0.1429}}
  >{\raggedright\arraybackslash}p{(\linewidth - 12\tabcolsep) * \real{0.1429}}
  >{\raggedright\arraybackslash}p{(\linewidth - 12\tabcolsep) * \real{0.1429}}
  >{\raggedright\arraybackslash}p{(\linewidth - 12\tabcolsep) * \real{0.1429}}@{}}
\toprule\noalign{}
\begin{minipage}[b]{\linewidth}\raggedright
Faixa
\end{minipage} & \begin{minipage}[b]{\linewidth}\raggedright
Normal R
\end{minipage} & \begin{minipage}[b]{\linewidth}\raggedright
Normal SPSS
\end{minipage} & \begin{minipage}[b]{\linewidth}\raggedright
Fórcipe R
\end{minipage} & \begin{minipage}[b]{\linewidth}\raggedright
Fórcipe SPSS
\end{minipage} & \begin{minipage}[b]{\linewidth}\raggedright
Cesárea R
\end{minipage} & \begin{minipage}[b]{\linewidth}\raggedright
Cesárea SPSS
\end{minipage} \\
\midrule\noalign{}
\endhead
\bottomrule\noalign{}
\endlastfoot
Precoces & 44,2\% & 45,4\% & 24,9\% & 24,5\% & 30,9\% & 30,1\% \\
Tardias & 40,2\% & 41,2\% & 27,6\% & 27,3\% & 32,2\% & 31,4\% \\
Adultas & 27,1\% & 27,5\% & 10,3\% & 10,2\% & 62,6\% & 62,3\% \\
\end{longtable}

\textbf{Problema:} As diferenças são pequenas (\textless{} 1,3 p.p.) mas
sistemáticas, resultado esperado de amostras ligeiramente distintas: R
aplica filtro de Robson e IG ≥ 22 semanas; SPSS usa apenas faixa etária
11--34 anos. Isso é documentado na nota de rodapé da Tabela 4 do docx.

\textbf{Correção sugerida:} Não é erro --- é diferença metodológica
documentada. A dissertação deve citar os valores do R (Tabela 4) para a
seção de resultados principal e, se mencionar os valores do SPSS,
indicar explicitamente que a amostra é diferente. A redação atual cita
os valores do R, que é a escolha correta.

\begin{center}\rule{0.5\linewidth}{0.5pt}\end{center}

\subsubsection{1.8 Tabela 8 --- cobertura muito baixa e desigual de
ind\_final para
adolescentes}\label{tabela-8-cobertura-muito-baixa-e-desigual-de-ind_final-para-adolescentes}

\textbf{Arquivo:} \texttt{tabelas\_dissertacao.docx} Tabela 8

\textbf{Problema:} Somando cesáreas + fórcipes do R (Tabela 4): precoces
= 300 partos operatórios, tardias = 496, adultas = 3.859. A Tabela 8 usa
denominadores de 138 (precoces), 227 (tardias) e 3.228 (adultas) ---
cobertura de apenas \textasciitilde46\% para adolescentes versus
\textasciitilde84\% para adultas. Isso significa que mais da metade dos
partos operatórios de adolescentes não tem indicação registrada na
variável \texttt{ind\_final} do SPSS.

\textbf{Problema derivado:} Os percentuais da Tabela 8 (ex: 34,8\%
sofrimento fetal nas precoces) representam apenas os 46\% com indicação
preenchida --- podem não ser representativos do total. A nota de rodapé
da tabela no docx menciona \texttt{ind\_final} mas não explicita essa
diferença de cobertura.

\textbf{Correção sugerida:} Acrescentar, no rodapé da Tabela 8, a
informação de cobertura: ``Denominadores refletem apenas os partos
operatórios com indicação registrada (precoces: 138/300 = 46\%; tardias:
227/496 = 46\%; adultas: 3.228/3.859 = 84\%). Resultado: as proporções
de adolescentes podem ser menos precisas que as das adultas.''
Considerar mencionar essa limitação também no texto de §4.4.

\begin{center}\rule{0.5\linewidth}{0.5pt}\end{center}

\subsubsection{1.9 Tabela 9 --- inteiramente ND (confirmado como
limitação)}\label{tabela-9-inteiramente-nd-confirmado-como-limitauxe7uxe3o}

\textbf{Arquivo:} \texttt{tabelas\_dissertacao.docx} Tabela 9;
\texttt{Relatorio\ Analise\ SPSS.md}

\textbf{Trecho relevante:} A variável \texttt{ind\_final} do SPSS agrega
indicações de cesárea e fórcipe sem estratificação por tipo de parto.

\textbf{Problema:} Não há como extrair indicações específicas de fórcipe
do SPSS sem cruzar \texttt{ind\_final} com \texttt{tipo\_parto} nos
microdados brutos. Isso é confirmado pela nota técnica no rodapé do
docx.

\textbf{Correção sugerida:} Isso é uma limitação genuína do dado, não um
erro. A Tabela 9 pode ser mantida com ND e nota explicativa, ou
substituída por texto descritivo em §4.4 reconhecendo a limitação. A
Tabela 8 (indicações combinadas de partos operatórios) já cobre o
assunto adequadamente. Se os microdados estiverem disponíveis, Eduardo
pode rodar o cruzamento \texttt{ind\_final\ ×\ tipo\_parto\ ==\ 3} no R
para extrair as indicações específicas de fórcipe.

\begin{center}\rule{0.5\linewidth}{0.5pt}\end{center}

\subsubsection{1.10 Figuras 5 e 6 --- citadas na redação mas arquivos
PNG
inexistentes}\label{figuras-5-e-6-citadas-na-redauxe7uxe3o-mas-arquivos-png-inexistentes}

\textbf{Arquivo:} \texttt{resultados\_redacao.docx} §4.4 vs
\texttt{results/figures/}

\textbf{Trecho:} \emph{``são apresentadas nas Tabelas 8 e 9, com
representação gráfica nas Figuras 5 e 6''}

\textbf{Problema:} Apenas cinco arquivos existem em
\texttt{results/figures/}: os quatro \texttt{fig\_obj2\_*} e
\texttt{fig\_obj5\_forest\_plot\_modelo4.png}. Não existe nenhum PNG
correspondente às Figuras 5 (indicações de cesárea) e 6 (indicações de
fórcipe).

\textbf{Correção sugerida:} Eduardo deve gerar as Figuras 5 e 6 rodando
o chunk correspondente do \texttt{index.qmd} e exportando com
\texttt{ggsave()}. Se as figuras de indicações de fórcipe ficarem todas
ND (por causa da Tabela 9), a Figura 6 deve ser suprimida e a referência
removida da redação §4.4.

\begin{center}\rule{0.5\linewidth}{0.5pt}\end{center}

\subsubsection{1.11 Figura 7 (forest plot) --- existe mas não é
referenciada na
redação}\label{figura-7-forest-plot-existe-mas-nuxe3o-uxe9-referenciada-na-redauxe7uxe3o}

\textbf{Arquivo:}
\texttt{results/figures/fig\_obj5\_forest\_plot\_modelo4.png} vs
\texttt{resultados\_redacao.docx} §4.5

\textbf{Problema:} O arquivo
\texttt{fig\_obj5\_forest\_plot\_modelo4.png} foi gerado e existe. O
§4.5 da redação não o referencia em nenhum momento.

\textbf{Correção sugerida:} Em §4.5, acrescentar uma referência à Figura
7 no parágrafo que introduz o Modelo 4. Exemplo: \emph{``Os resultados
do Modelo 4 estão apresentados na Tabela 10 e ilustrados graficamente na
Figura 7.''} Ajustar a numeração das figuras no docx principal conforme
o plano (Figura 7 = forest plot).

\begin{center}\rule{0.5\linewidth}{0.5pt}\end{center}

\subsubsection{1.12 Modelo 4 --- verificação dos OR e IC 95\%
(resultado: SEM
ERROS)}\label{modelo-4-verificauxe7uxe3o-dos-or-e-ic-95-resultado-sem-erros}

\textbf{Arquivos:} \texttt{tab10\_modelo4\_spss.csv},
\texttt{tabelas\_dissertacao.docx} Tabela 10,
\texttt{Relatorio\ Analise\ SPSS.md} linhas 123--135

Verificação completa de todos os 11 coeficientes:

\begin{longtable}[]{@{}
  >{\raggedright\arraybackslash}p{(\linewidth - 6\tabcolsep) * \real{0.2500}}
  >{\raggedright\arraybackslash}p{(\linewidth - 6\tabcolsep) * \real{0.2500}}
  >{\raggedright\arraybackslash}p{(\linewidth - 6\tabcolsep) * \real{0.2500}}
  >{\raggedright\arraybackslash}p{(\linewidth - 6\tabcolsep) * \real{0.2500}}@{}}
\toprule\noalign{}
\begin{minipage}[b]{\linewidth}\raggedright
Variável
\end{minipage} & \begin{minipage}[b]{\linewidth}\raggedright
SPSS (Relatorio)
\end{minipage} & \begin{minipage}[b]{\linewidth}\raggedright
Tabela 10 docx
\end{minipage} & \begin{minipage}[b]{\linewidth}\raggedright
Situação
\end{minipage} \\
\midrule\noalign{}
\endhead
\bottomrule\noalign{}
\endlastfoot
Faixa etária adulta & OR=1,79 {[}1,12--2,86{]} p=0,015 & OR=1,79
{[}1,12--2,86{]} p=0,015 & ✓ \\
DHEG & OR=0,40 {[}0,32--0,50{]} p\textless0,001 & OR=0,40
{[}0,32--0,50{]} p\textless0,001 & ✓ \\
Robson 2 & OR=4,38 {[}3,30--5,82{]} p\textless0,001 & OR=4,38
{[}3,30--5,82{]} p\textless0,001 & ✓ \\
Robson 3 & OR=0,48 {[}0,35--0,66{]} p\textless0,001 & OR=0,48
{[}0,35--0,66{]} p\textless0,001 & ✓ \\
Robson 4 & OR=1,86 {[}1,33--2,59{]} p\textless0,001 & OR=1,86
{[}1,33--2,59{]} p\textless0,001 & ✓ \\
Robson 5 & OR=13,65 {[}9,88--18,85{]} p\textless0,001 & OR=13,65
{[}9,88--18,85{]} p\textless0,001 & ✓ \\
Robson 6 & OR=23,65 {[}5,49--101,89{]} p\textless0,001 & OR=23,65
{[}5,49--101,89{]} p\textless0,001 & ✓ \\
Robson 7 & OR=22,86 {[}13,05--40,08{]} p\textless0,001 & OR=22,86
{[}13,05--40,08{]} p\textless0,001 & ✓ \\
Robson 9 & OR=27,66 {[}6,49--117,89{]} p\textless0,001 & OR=27,66
{[}6,49--117,89{]} p\textless0,001 & ✓ \\
Robson 10 & OR=2,69 {[}2,05--3,52{]} p\textless0,001 & OR=2,69
{[}2,05--3,52{]} p\textless0,001 & ✓ \\
Constante & −0,384 / p=0,073 / OR=0,68 & −0,384 / p=0,073 / OR=0,68 &
✓ \\
\end{longtable}

\textbf{Conclusão:} Tabela 10 está perfeita. Sem correções necessárias.

\begin{center}\rule{0.5\linewidth}{0.5pt}\end{center}

\subsubsection{1.13 Tabela 7 --- arredondamento do p-valor do Grupo
3}\label{tabela-7-arredondamento-do-p-valor-do-grupo-3}

\textbf{Arquivo:} \texttt{tab\_obj2\_testes\_por\_robson.csv} vs
\texttt{tabelas\_dissertacao.docx} Tabela 7 vs
\texttt{resultados\_redacao.docx} §4.3.2

\begin{longtable}[]{@{}ll@{}}
\toprule\noalign{}
Fonte & Grupo 3 p-valor \\
\midrule\noalign{}
\endhead
\bottomrule\noalign{}
\endlastfoot
CSV primário & p = 0,0014 \\
Tabela 7 docx & p = 0,001 \\
Redação §4.3.2 & p = 0,001 \\
\end{longtable}

\textbf{Problema:} O CSV primário mostra p = 0,0014. O docx e a redação
arredondam para p = 0,001. Não é erro, mas não é o valor exato. Segundo
os critérios de elegibilidade estatística da dissertação (α = 5\%),
ambos estão acima de p = 0,001.

\textbf{Correção sugerida:} Padronizar para p = 0,001 ou reportar o
valor exato p = 0,0014 em ambos. Para revistas científicas e banca de
mestrado, reportar p = 0,001 é aceitável (arredondamento a 3 casas).
Manter a consistência entre docx e redação --- ambos já usam 0,001,
então está OK.

\begin{center}\rule{0.5\linewidth}{0.5pt}\end{center}

\subsection{2. Itens em Branco (ND e
Placeholders)}\label{itens-em-branco-nd-e-placeholders}

\subsubsection{2.1 Cinco placeholders na redação --- ação
obrigatória}\label{cinco-placeholders-na-redauxe7uxe3o-auxe7uxe3o-obrigatuxf3ria}

\textbf{Arquivo:} \texttt{resultados\_redacao.docx} §4.1 (parágrafo 2 e
parágrafo de §4.1.1)

\begin{longtable}[]{@{}
  >{\raggedright\arraybackslash}p{(\linewidth - 6\tabcolsep) * \real{0.2500}}
  >{\raggedright\arraybackslash}p{(\linewidth - 6\tabcolsep) * \real{0.2500}}
  >{\raggedright\arraybackslash}p{(\linewidth - 6\tabcolsep) * \real{0.2500}}
  >{\raggedright\arraybackslash}p{(\linewidth - 6\tabcolsep) * \real{0.2500}}@{}}
\toprule\noalign{}
\begin{minipage}[b]{\linewidth}\raggedright
Placeholder
\end{minipage} & \begin{minipage}[b]{\linewidth}\raggedright
Descrição
\end{minipage} & \begin{minipage}[b]{\linewidth}\raggedright
Fonte recomendada
\end{minipage} & \begin{minipage}[b]{\linewidth}\raggedright
Valor provisório (R)
\end{minipage} \\
\midrule\noalign{}
\endhead
\bottomrule\noalign{}
\endlastfoot
\texttt{{[}N\_TOTAL{]}} & Total de partos analisados & Tabela 4 rodapé
--- \texttt{tab\_obj2\_taxa\_cesarea.csv} (soma de todos os grupos) &
\textbf{6.663} \\
\texttt{{[}N\_ADOL{]}} & Total de adolescentes & Soma precoces + tardias
da mesma fonte & \textbf{1.368} (538+830) \\
\texttt{{[}N\_PREC{]}} & Adolescentes precoces (11--15) &
\texttt{tab\_obj2\_taxa\_cesarea.csv} soma Grupo 1-10 para precoces &
\textbf{538} \\
\texttt{{[}N\_TARD{]}} & Adolescentes tardias (16--19) &
\texttt{tab\_obj2\_taxa\_cesarea.csv} soma Grupo 1-10 para tardias &
\textbf{830} \\
\texttt{{[}N\_ADUL{]}} & Adultas (20--34) &
\texttt{tab\_obj2\_taxa\_cesarea.csv} soma Grupo 1-10 para adultas &
\textbf{5.295} \\
\end{longtable}

\begin{quote}
\textbf{Atenção:} Os valores acima são os valores do docx de tabelas
(baseados em \texttt{tab\_obj2\_taxa\_cesarea.csv}). Se Eduardo
revalidar o R e os dois CSVs primários convergirem para 6.646 (com
tardias=829 e adultas=5.279), use esses valores menores. A decisão de
qual N usar deve vir ANTES de colar qualquer número no docx principal.
\end{quote}

\subsubsection{2.2 Tabela 9 --- todas as células
ND}\label{tabela-9-todas-as-cuxe9lulas-nd}

\textbf{Arquivo:} \texttt{tabelas\_dissertacao.docx} Tabela 9

\textbf{Ação recomendada:} Duas opções: - \textbf{Opção A
(recomendada):} Eduardo cruza \texttt{ind\_final\ ×\ tipo\_parto\ ==\ 3}
nos microdados brutos do R, extrai as indicações específicas de fórcipe
e preenche a Tabela 9. - \textbf{Opção B:} Suprimir a Tabela 9,
transformar §4.4 ``Indicações de fórcipe'' em um parágrafo descritivo
que reconhece a limitação do dado e cita apenas a frequência global de
fórcipe (já disponível na Tabela 4).

\subsubsection{2.3 Tabela 11 (desfechos neonatais) --- CSV existe, não
está no
docx}\label{tabela-11-desfechos-neonatais-csv-existe-nuxe3o-estuxe1-no-docx}

\textbf{Arquivo:} \texttt{tab11\_desfechos\_neonatais\_spss.csv} existe;
\texttt{tabelas\_dissertacao.docx} não contém Tabela 11.

\textbf{Ação recomendada:} Decisão pendente da Letícia (ver PLANO §3.6).
Se desfechos neonatais forem incluídos como objetivo secundário, a
Tabela 11 deve ser formatada e inserida no docx. O CSV de base existe,
com quatro colunas de texto narrativo (Apgar e óbito fetal não têm
valores numéricos precisos para tardias). Se for excluída da seção
Resultados, levar os principais achados para a Discussão como
``elementos contextuais''.

\begin{center}\rule{0.5\linewidth}{0.5pt}\end{center}

\subsection{3. Recomendações Finais}\label{recomendauxe7uxf5es-finais}

\subsubsection{Para a Letícia}\label{para-a-letuxedcia}

\begin{enumerate}
\def\labelenumi{\arabic{enumi}.}
\item
  \textbf{Decisão prioritária --- o N:} Antes de colar qualquer tabela
  ou parágrafo no docx principal, definir com Eduardo qual N usar para
  as tabelas de análise (Tabelas 4--7): 6.646 (CSV primário de vias) ou
  6.663 (CSV primário de taxa/Robson). Essa diferença de 17 casos
  precisa ser investigada e resolvida, não contornada. Uma vez definido,
  preencher os cinco placeholders.
\item
  \textbf{Substituições no texto original da dissertação:}

  \begin{itemize}
  \tightlist
  \item
    §154: substituir ``8.300 participantes, sendo 1.348 parturientes
    adolescentes e 6.952 adultas'' pelos N definitivos.
  \item
    Legendas das Tabelas 1 e 2 originais: substituir ``1.677
    adolescentes e 7.047 adultas'' pelos N definitivos.
  \item
    §160: substituir ``p valor neste caso é de 0,06'' por ``p = 0,337''.
  \end{itemize}
\item
  \textbf{Decisão sobre a Tabela 9 e os desfechos neonatais:} São duas
  decisões independentes que determinam se a dissertação terá 10 ou 11
  tabelas. Recomendo tomar essas decisões antes de montar o docx final.
\item
  \textbf{Verificar o valor de 56,3\% (cor/raça):} Confirmar se o SPSS
  reporta exatamente 56,3\% para o grupo adolescente combinado ou se o
  valor preciso é diferente. Ajustar §4.1.2 conforme necessário.
\item
  \textbf{Estilo --- itálico para \emph{p}:} Ao colar os parágrafos da
  nova redação no docx principal, aplicar itálico a toda ocorrência de
  \emph{p} estatístico nos textos corridos (ex: \emph{p} \textless{}
  0,001; \emph{p} = 0,337). Verifique também que os números nas tabelas
  usam vírgula decimal de forma consistente --- as tabelas do docx já
  estão corretas, mas ao formatar no Word pode ser que formatações
  automáticas convertam vírgula para ponto.
\end{enumerate}

\subsubsection{Para o Eduardo}\label{para-o-eduardo}

\begin{enumerate}
\def\labelenumi{\arabic{enumi}.}
\item
  \textbf{Causa raiz da divergência de N (6.646 vs 6.663):} Rerodar
  \texttt{quarto\ render\ index.qmd} e verificar se
  \texttt{tab\_obj2\_vias\_parto\_geral.csv} e
  \texttt{tab\_obj2\_taxa\_cesarea.csv} produzem somas concordantes. Se
  divergirem, identificar o chunk com o filtro diferente. A hipótese
  mais provável é que um dos CSVs usa \texttt{dados\_analise} e o outro
  usa \texttt{dados\_s4} (antes de criar \texttt{faixa\_etaria}),
  capturando casos com idade fora do range.
\item
  \textbf{Deletar ou arquivar CSVs derivados conflitantes:} Os arquivos
  \texttt{results/tabelas\_dissertacao/tab05\_robson\_faixa\_via.csv} e
  \texttt{tab06\_taxa\_cesarea\_robson.csv} têm valores diferentes dos
  primários e podem causar confusão. Arquivar em subpasta
  \texttt{\_deprecated/} ou regenerar do zero garantindo que usam a
  mesma base.
\item
  \textbf{Gerar Figuras 5 e 6:} Rodar os chunks de indicações do
  \texttt{index.qmd} e exportar como PNG em \texttt{results/figures/}
  com nomenclatura \texttt{fig\_obj3\_indicacoes\_cesarea.png} e
  \texttt{fig\_obj3\_indicacoes\_forcipe.png} (ou similar). Se a Figura
  6 depender de dados que resultam todos em ND, suprimi-la.
\item
  \textbf{Adicionar Figura 7 na redação:} Em
  \texttt{resultados\_redacao.docx} §4.5, inserir referência à Figura 7
  (forest plot do Modelo 4). O arquivo já existe em
  \texttt{results/figures/fig\_obj5\_forest\_plot\_modelo4.png}.
\item
  \textbf{Tabela 9 (se Opção A):} Cruzar
  \texttt{ind\_final\ ×\ tipo\_parto} nos microdados para extrair
  indicações específicas de fórcipe. Isso resolve tanto a Tabela 9
  quanto o detalhe de baixa cobertura da Tabela 8 para adolescentes.
\item
  \textbf{Nota de rodapé na Tabela 8:} Acrescentar a informação de
  cobertura: ``Partos operatórios com indicação registrada: 138/300
  (46\%) precoces; 227/496 (46\%) tardias; 3.228/3.859 (84\%) adultas.''
\end{enumerate}

\begin{center}\rule{0.5\linewidth}{0.5pt}\end{center}

\subsection{Apêndice --- Rastreabilidade dos
Placeholders}\label{apuxeandice-rastreabilidade-dos-placeholders}

\begin{longtable}[]{@{}llll@{}}
\toprule\noalign{}
Placeholder & Valor R (6.663) & Valor R (6.646) & Tabela de origem \\
\midrule\noalign{}
\endhead
\bottomrule\noalign{}
\endlastfoot
{[}N\_TOTAL{]} & 6.663 & 6.646 & Rodapé Tabela 4 \\
{[}N\_ADOL{]} & 1.368 & 1.367 & Soma precoces + tardias \\
{[}N\_PREC{]} & 538 & 538 & Tabela 4, coluna precoces \\
{[}N\_TARD{]} & 830 & 829 & Tabela 4, coluna tardias \\
{[}N\_ADUL{]} & 5.295 & 5.279 & Tabela 4, coluna adultas \\
\end{longtable}

\begin{center}\rule{0.5\linewidth}{0.5pt}\end{center}

\emph{Relatório gerado automaticamente por revisão cruzada de 13
arquivos. Última atualização: 19/05/2026.}

\end{document}
